\section{渊源和效力}
\begin{enumerate}
  \item 渊源，有且只有三项。
        \begin{enumerate}
          \item \zd{条约}。核心区别范围不一样，约束\zd{缔约国}。
          \item \zd{习惯}。\zd{实践产物}；约束\zd{所有}，\zd{持续反对者（核相关）除外}；证据来源\zd{官方}；不成文。
          \item \zd{一般法律原则}。太多了。
        \end{enumerate}
  \item 国际法的基本原则。适用国际法所有领域的。具有\zd{强行法}性质。
        \begin{enumerate}
          \item 国家主权平等。
          \item 不干涉内政（国家管辖权概念）。
          \item 不得使用威胁 \hz 武力原则。并\zd{不是禁止一切}武力。自卫 + 安理会授权。
          \item 和平（\qd{\bdy 非武力}）解决国际争端。\qd{自卫是武力，但仍然称之为和平手段}。
          \item 民族自决。所有。独立权只适用于\zd{殖民统治}下。
          \item 善意履行国际义务。
        \end{enumerate}
\end{enumerate}
